\chapter{Scope and status}

\section{Conventions and elementary objects}

We work in the complex plane $\C$ with its Euclidean metric and write
\[
\overline B(a,r):=\{z\in\C:|z-a|\le r\},\qquad
B(a,r):=\{z\in\C:|z-a|<r\}.
\]
The natural numbers are $\N=\{0,1,2,\ldots\}$.  For a map $g$, the notation
$g^{\circ n}$ means the $n$-fold iterate, with $g^{\circ0}=\id$.

\begin{definition}[Quadratic family and critical orbit]
\label{def:quadratic-family}
For $c\in\C$ let
\[
 f_c(z):=z^2+c,\qquad
 p_n(c):=f_c^{\circ n}(0).
\]
The orbit of a point $z$ is
\[
 \operatorname{Orb}_c(z):=\{f_c^{\circ n}(z):n\in\N\}.
\]
\end{definition}

\begin{definition}[Filled Julia set and Julia set]
\label{def:filled-julia}
The filled Julia set and Julia set of $f_c$ are
\[
 \Kc:=\{z\in\C:\sup_{n\in\N}|f_c^{\circ n}(z)|<\infty\},
 \qquad
 \Jc:=\partial\Kc.
\]
\end{definition}

\begin{definition}[Mandelbrot set]
\label{def:mandelbrot}
The Mandelbrot set is the critical-orbit locus
\[
 \M:=\{c\in\C:0\in\Kc\}
   =\{c\in\C:\sup_{n\in\N}|p_n(c)|<\infty\}.
\]
\end{definition}

\begin{theorem}[Escape-radius criterion]
\label{thm:escape-radius}
If $|z|>\max\{2,|c|\}$, then
\[
 |f_c^{\circ n}(z)|\longrightarrow\infty.
\]
Consequently, for $c\in\C$,
\[
 c\in\M\quad\Longleftrightarrow\quad
 |p_n(c)|\le2\text{ for every }n\in\N,
\]
and $\M\subseteq\overline B(0,2)$.
\end{theorem}

\begin{proof}
If $|z|>\max\{2,|c|\}$, then
\[
 |f_c(z)|\ge |z|^2-|c|>|z|^2-|z|>|z|.
\]
The same inequality applies inductively to every subsequent iterate.  The
critical orbit statement follows by applying the inequality to the first
iterate which leaves $\overline B(0,2)$; the converse is immediate from the
definition of $\Kc$.
\end{proof}

\begin{definition}[Finite outer approximations]
\label{def:finite-outer}
For $N\in\N$ set
\[
 O_N:=\left\{c\in\C:
 |c|\le2,\quad |p_j(c)|\le2\text{ for every }j\le N\right\}.
\]
\end{definition}

\begin{theorem}[Compact nested approximation]
\label{thm:finite-outer}
The sets $O_N$ are compact, $O_{N+1}\subseteq O_N$, and
\[
 \M\subseteq O_N,\qquad
 \bigcap_{N\in\N}O_N=\M.
\]
\end{theorem}

\begin{proof}
Each $p_j$ is a polynomial in $c$, so the defining inequalities are closed.
The escape-radius criterion gives $O_N\subseteq\overline B(0,2)$, hence
compactness.  The nesting and inclusion of $\M$ are immediate, and the
intersection equality follows by checking the $N$-th orbit inequality at
the stage $N$.
\end{proof}

\section{Local connectivity and components}

\begin{definition}[Local connectedness at a point]
\label{def:local-connectedness}
Let $X$ be a topological space and $x\in X$.  We say that $X$ is locally
connected at $x$ if every neighbourhood $U\in\mathcal N_X(x)$ contains a
neighbourhood $V\in\mathcal N_X(x)$ such that $V$ is preconnected and
$V\subseteq U$.  The space is locally connected if it is locally connected
at every point.
\end{definition}

\begin{definition}[Component modulus]
\label{def:component-modulus}
For $x,y\in\M$ define
\[
 \omega_\M(r):=
 \sup_{\substack{x,y\in\M\\|x-y|<r}}
 \inf\{\diam C:C\subseteq\M\text{ compact and connected},\
 x,y\in C\}.
\]
The convention is $\inf\varnothing=+\infty$.
\end{definition}

\begin{theorem}[Metric criterion for local connectivity]
\label{thm:metric-lc}
For a compact metric space $X$, local connectedness is equivalent to
\[
 \forall\varepsilon>0\ \exists\delta>0\ \forall x,y\in X,\qquad
 d(x,y)<\delta\Longrightarrow
 \exists C\subseteq X:
 \begin{cases}
 C\text{ is compact and connected},\\
 x,y\in C,\qquad \diam C<\varepsilon.
 \end{cases}
\]
In particular,
\[
 \M\text{ is locally connected}
\quad\Longleftrightarrow\quad
 \lim_{r\downarrow0}\omega_\M(r)=0.
\]
\end{theorem}

\section{Green functions and external coordinates}

\begin{definition}[Escape-rate Green function]
\label{def:green-function}
For $c,z\in\C$ define
\[
 G_c(z):=\lim_{n\to\infty}2^{-n}\log^+
       |f_c^{\circ n}(z)|,
 \qquad \log^+ t:=\max\{\log t,0\}.
\]
\end{definition}

\begin{theorem}[Green-function properties]
\label{thm:green-properties}
For every $c\in\C$, the limit in Definition~\ref{def:green-function}
exists and $G_c$ is continuous.  Moreover,
\[
 G_c(f_c(z))=2G_c(z),\qquad
 G_c(z)=0\Longleftrightarrow z\in\Kc,
\]
and $G_c$ is harmonic on $\C\setminus\Kc$.  On the basin of infinity there
is a unique normalized Böttcher coordinate $\phi_c$ satisfying
\[
 \phi_c(f_c(z))=\phi_c(z)^2,\qquad
 \phi_c(z)=z+O(1)\quad(z\to\infty),
\]
with $G_c(z)=\log|\phi_c(z)|$.
\end{theorem}

\begin{definition}[External ray]
\label{def:external-ray}
For $\theta\in\R/\Z$, the external ray of angle $\theta$ is
\[
 R_c(\theta):=
 \phi_c^{-1}\{re^{2\pi i\theta}:r>1\}.
\]
Landing of $R_c(\theta)$ means that its closure in $\C$ contains a unique
finite endpoint.
\end{definition}

\begin{conditionalstep}
The Böttcher coordinate and the Green function describe the basin of
infinity.  To obtain a parameter-plane theorem one additionally needs a
joint phase--parameter construction: rays and equipotentials must move with
$c$, their finite graphs must cut both dynamical and parameter planes, and
the resulting parapuzzle pieces must have the same combinatorics.  This is
the first genuinely non-elementary bridge in the MLC program.
\end{conditionalstep}

\section{Puzzles, parapuzzles, and modulus}

\begin{definition}[Puzzle graph and puzzle piece]
\label{def:puzzle-piece}
Fix a finite collection of landing external rays and equipotential arcs
which is forward invariant under angle doubling.  The resulting graph at
depth $n$ is denoted by $\Gamma_n(c)$.  A dynamical puzzle piece is a
connected component
\[
 P_n(c,z):=\operatorname{Comp}_z(\C\setminus\Gamma_n(c))
\]
containing $z$.  The critical puzzle piece is $P_n(c,0)$.
\end{definition}

\begin{definition}[Parapuzzle piece]
\label{def:parapuzzle-piece}
Let $\mathcal G_n$ be the corresponding parameter graph obtained from the
phase--parameter correspondence.  A parameter parapuzzle piece is a
connected component of $\C\setminus\mathcal G_n$ with fixed finite
combinatorial data.  We write $\mathcal P_n(c)$ for the piece containing
$c$.
\end{definition}

\begin{definition}[Annulus modulus]
\label{def:annulus-modulus}
If $A=U\setminus\overline V$ is a doubly connected domain conformally
equivalent to $\{1<|z|<R\}$, define
\[
 \modulus(A):=\frac{1}{2\pi}\log R.
\]
\end{definition}

\begin{theorem}[Grötzsch--Yoccoz shrinking criterion]
\label{thm:yoccoz-shrink}
Suppose that $P_{n+1}(c,0)\subseteq P_n(c,0)$ are connected puzzle pieces
and that the annuli
\[
 A_n:=P_n(c,0)\setminus\overline{P_{n+1}(c,0)}
\]
are nondegenerate.  If
\[
 \sum_{n\in\N}\modulus(A_n)=+\infty,
\]
then
\[
 \bigcap_{n\in\N}\overline{P_n(c,0)}=\{0\}.
\]
\end{theorem}

\begin{theorem}[Parameter shrinking implication]
\label{thm:parameter-shrinking}
Assume that the phase--parameter correspondence sends the connected
dynamical pieces in Theorem~\ref{thm:yoccoz-shrink} to connected
parapuzzle pieces $\mathcal P_n(c)$ and that
\[
 \bigcap_{n\in\N}\overline{\mathcal P_n(c)}=\{c\}.
\]
Then $\M$ is locally connected at $c$.
\end{theorem}

\begin{conditionalstep}
Theorem~\ref{thm:yoccoz-shrink} is a dynamical-plane modulus statement.
Theorem~\ref{thm:parameter-shrinking} is the parameter-plane conclusion.
The latter is not obtained by replacing a parapuzzle by a translated full
Green sublevel; the graph-cut, landing, and phase--parameter hypotheses have
to be formalized separately.
\end{conditionalstep}

\section{The target}

Let
\[
 f_c(z)=z^2+c,\qquad p_0(c)=0,\qquad p_{n+1}(c)=f_c(p_n(c)),
\]
and let
\[
 \M=\{c\in\C:\sup_n |p_n(c)|<\infty\}
\]
be the Mandelbrot set.  The global MLC conjecture is
\[
 \mathsf{MLC}:\qquad \LocallyConnectedSpace(\M).
\]
The pointwise statement is weaker in presentation but is the form used by
the literature: for $c\in\M$, every neighbourhood of $c$ in $\M$ contains
a connected neighbourhood of $c$.

The first distinction of this review is therefore:
\begin{enumerate}
\item a theorem about a parameter class;
\item a theorem about local connectivity of the corresponding Julia set;
\item a theorem about local connectivity of the Mandelbrot set at those
      parameters; and
\item the global statement $\mathsf{MLC}$.
\end{enumerate}
These statements are related by the phase--parameter machinery, but they are
not interchangeable.  The Lean declarations
\lean{Mlc.Literature.MandelbrotLocallyConnectedAt} and
\lean{Mlc.Literature.GlobalMLC} keep this distinction explicit.

\section{A status picture}

\clearpage
\begin{landscape}
\thispagestyle{plain}
\begin{center}
\includegraphics[width=0.82\linewidth]{docs/figures/mlc_status_render.pdf}
\captionof{figure}{A numerical and literature-status picture of the Mandelbrot set.  It is not a
numerically certified characteristic function of the unresolved set.  The
high-depth render, zoom panels, and component--theorem ledger compare the
geometry of $M$ with the current proof inventory.  The coloured labels are
literature classes, not claims that a whole visible boundary arc has already
been proved locally connected.}
\label{fig:mlc-status}
\end{center}
\end{landscape}
\clearpage


The green interior in Figure~\ref{fig:mlc-status} is not the interesting
part of the conjecture: open subsets of the plane are locally connected.  At
the boundary, the established regions are best recorded as classes of
parameters rather than as literal open subsets of the drawing.  In
particular, ``bounded type'' and ``Feigenbaum'' refer to renormalization
conditions, not to a visually isolated bulb.

\section{What remains open}

The full conjecture remains open.  The main unresolved frontier is the
parameter-plane control needed for general infinitely renormalizable
quadratics, in particular cases with unbounded combinatorics.  Recent
bounded-type satellite and Feigenbaum results enlarge the proved set of
parameters, but do not supply a uniform theorem for every boundary
parameter.

The repository must consequently distinguish:
\[
\begin{array}{c}
\text{proved Lean theorem}\\
\Downarrow\\
\text{formal theorem interface with explicit hypotheses}\\
\Downarrow\\
\text{external analytic result still to be formalized}\\
\Downarrow\\
\mathsf{MLC}.
\end{array}
\]
Replacing an external result by an `axiom` would make the final theorem
shorter, but would not constitute the requested formalization.
